%************************************************
\chapter{Variational Autoencoders}\label{ch:vaes}
%************************************************

The core initial idea for the understanding of the Variational Autoencoders (VAEs) is the concept of \textbf{Autoencoders} (AEs). 
This are a type of model that contains two elements, the encoder and the decoder. The encoder is a function that maps the input data 
$\mathbf{x}\in\mathbb{R}^d$ to a low-dimensional representation $\mathbf{z}\in\mathbb{R}^k$, which is usually much more smaller than
the input dimension, $k<<d$. This representations live in a space with a more compact and interpretable structure, called the latent space. 
Then, the decoder is a function that maps the latent representation $\mathbf{z}$ back to the input space, trying to reconstruct the 
original data point $\mathbf{x}$. \\
Once we have an autoencoder model trained, a natural question is how to use it as a generative model. One idea is to randomly sample
a point $\mathbf{z}$ from the latent space and then pass it through the decoder to get a synthetic data point $\mathbf{x'}$. This method will 
not produce realistic samples because the latent space is disorganized and irregular. Another option is to encode and image and sample 
neighboring points in the latent space that could generate similar samples, but the reality is that the output are not meaningful samples. This 
shows how poorly structured the latent space is. The aim of VAEs is to transform the latent space of an autoencoder into a well-structured space 
that follows a prior distribution, which allows us to sample from it and generate realistic samples. \\
\\
It is also important to introduce the concept of latent variable models. This type of variables are part of the model 
but are not observed in the data, usually denoted as $\mathbf{z}$, which are the hidden factors of variation. The goal of latent variable models is to 
learn the joint distribution $p_{\theta}(\mathbf{x}, \mathbf{z})$ of the observed data $\mathbf{x}$ and the latent variables $\mathbf{z}$. 
The marginal likelihood, model evidence or marginal distribution over the observed data is given by:
\begin{equation}\label{eq:p_x}
    p_{\theta}(\mathbf{x}) = \int p_{\theta}(\mathbf{x}, \mathbf{z}) d\mathbf{z}.
\end{equation}
When we use neural networks to parameterize the latent variable models we call them \textbf{Deep Latent Variable Models} (DLVMs). 
The simple and most common DLVM is given by the following factorization:
\begin{equation}
    p_{\theta}(\mathbf{x}, \mathbf{z}) = p(\mathbf{z}) p_{\theta}(\mathbf{x} | \mathbf{z}),
\end{equation}
where $p(\mathbf{z})$ is the \textit{prior distribution} over the latent variables and $p_{\theta}(\mathbf{x} | \mathbf{z})$ 
is the likelihood of the data given the latent variables, also called the \textit{generator distribution}. The joint distribution $p_{\theta}(\mathbf{x}, \mathbf{z})$ 
is related through the Bayes' rule to the following posterior distribution:
\begin{equation}
    p_{\theta}(\mathbf{z} | \mathbf{x}) = \frac{p_{\theta}(\mathbf{x} | \mathbf{z}) p(\mathbf{z})}{p_{\theta}(\mathbf{x})},
\end{equation}
where we have an intractability problem due to the marginal likelihood $p_{\theta}(\mathbf{x})$ which is given by an integral over the latent variables (Equation \ref{eq:p_x}). 
To overcome this problem we can use variational inference (VI) which is a technique to approximate the posterior distribution $p_{\theta}(\mathbf{z} | \mathbf{x})$ 
with a simpler distribution $q_{\phi}(\mathbf{z} | \mathbf{x})$ that is parameterized by $\phi$.  \\
\\
This is the main idea behind \textbf{Variational Autoencoders} (VAEs) \cite{kingma2022autoencodingvariationalbayes} which are a type of DLVMs that use VI to learn the latent variable model. It defines a 
parametric inference model $q_{\phi}(\mathbf{z} | \mathbf{x})$, also called an \textit{encoder}, where $\phi$ are the parameters. The idea is to optimize the parameters such that: 
\begin{equation}
    q_{\phi}(\mathbf{z} | \mathbf{x}) \approx p_{\theta}(\mathbf{z} | \mathbf{x}).
\end{equation}
This VAE formulation allows us to transform an autoencoder model with an unstructured latent space into a generative model with a latent space that follows a prior 
distribution, typically a Gaussian distribution $p(\mathbf{z}) = \mathcal{N}(\mathbf{z}; \mathbf{0}, \mathbf{I})$, which is easy to sample from. \\

\begin{figure}
\centering
\begin{tikzpicture}[
    >=Stealth,
    node distance=1.0cm, % Very tight spacing
    % Styles for nodes
    block/.style={draw=blue!70!black, fill=blue!5, rectangle, minimum height=1cm, minimum width=2.2cm, align=center, rounded corners=2pt, thick, font=\small\bfseries},
    latentblock/.style={draw=green!70!black, fill=green!5, circle, minimum size=0.8cm, align=center, thick, font=\small},
    io/.style={draw=black!80, fill=white, circle, inner sep=1pt, minimum size=0.8cm, align=center, thick, font=\small}
]

    % Nodes with labels placed below them compactly
    \node (input) [io, label=below:{\scriptsize Input}] {$x$};
    
    \node (encoder) [block, right=of input] {Encoder \\[-0.5ex] \scriptsize $q_{\phi}(z|x)$};
    
    \node (z) [latentblock, right=of encoder, label=below:{\scriptsize Latent}] {$z$};
    
    \node (decoder) [block, right=of z] {Decoder \\[-0.5ex] \scriptsize $p_{\theta}(x|z)$};
    
    \node (output) [io, right=of decoder, label=below:{\scriptsize Recon}] {$x'$};

    % Compact Arrows
    \draw [->, thick, draw=black!70] (input) -- (encoder);
    \draw [->, thick, draw=black!70] (encoder) -- (z);
    \draw [->, thick, draw=black!70] (z) -- (decoder);
    \draw [->, thick, draw=black!70] (decoder) -- (output);

    \path ([yshift=0.8cm]current bounding box.north);

\end{tikzpicture}
\caption{Simplified Variational Autoencoder (VAE) architecture mapping an input $x$ to a latent representation $z$, and decoding it to a reconstruction $x'$.}\label{diag:vae}
\end{figure}


% ******
% Training 
% ******
\section{Training of VAEs} 
For the training of the VAEs we can not directly optimize the marginal likelihood $p_{\theta}(\mathbf{x})$ due to the intractability problem, but we can optimize a lower bound of the marginal likelihood
called the \textbf{Evidence Lower Bound} (ELBO). Let's derive the log-likelihood of the data point $\mathbf{x}$ and see how we can get the ELBO using the Jensen's inequality \cite{Kingma_2019}:
\begin{align}
    \log p_{\theta}(\mathbf{x}) &= \log\int p_{\theta}(\mathbf{x},\mathbf{z})d\mathbf{z} \\ %= \mathbb{E}_{\mathbf{z}\sim q_{\phi}(\mathbf{z} | \mathbf{x})} \left[\log p_{\theta}(\mathbf{x})\right] \\
    &= \log\int q_{\phi}(\mathbf{z}|\mathbf{x})\frac{p_{\theta}(\mathbf{x}, \mathbf{z})}{q_{\phi}(\mathbf{z}|\mathbf{x})}d\mathbf{z} \\
    &= \log \mathbb{E}_{\mathbf{z}\sim q_{\phi}(\mathbf{z}|\mathbf{x})} \left[\frac{p_{\theta}(\mathbf{x},\mathbf{z})}{q_{\phi}(\mathbf{z}|\mathbf{x})}\right] \\
    &\geq \mathbb{E}_{\mathbf{z}\sim q_{\phi}(\mathbf{z}|\mathbf{x})} \left[\log \frac{p_{\theta}(\mathbf{x},\mathbf{z})}{q_{\phi}(\mathbf{z}|\mathbf{x})}\right] \\
    &= \mathbb{E}_{\mathbf{z}\sim q_{\phi}(\mathbf{z}|\mathbf{x})} \left[\log \frac{p_{\theta}(\mathbf{x}|\mathbf{z})p(\mathbf{z})}{q_{\phi}(\mathbf{z}|\mathbf{x})}\right] \\
    &= \mathbb{E}_{\mathbf{z}\sim q_{\phi}(\mathbf{z}|\mathbf{x})} \left[\log \frac{p_{\theta}(\mathbf{x}|\mathbf{z})}{\frac{q_{\phi}(\mathbf{z}|\mathbf{x})}{p(\mathbf{z})}}\right] \\
    &= \mathbb{E}_{\mathbf{z}\sim q_{\phi}(\mathbf{z}|\mathbf{x})} \left[\log p_{\theta}(\mathbf{x}|\mathbf{z}) - \log \frac{q_{\phi}(\mathbf{z}|\mathbf{x})}{p(\mathbf{z})}\right] \\
    &= \underbrace{\mathbb{E}_{\mathbf{z}\sim q_{\phi}(\mathbf{z}|\mathbf{x})} \left[\log p_{\theta}(\mathbf{x}|\mathbf{z})\right]}_{\text{Reconstruction}} - \underbrace{\mathbb{E}_{\mathbf{z}\sim q_{\phi}(\mathbf{z}|\mathbf{x})} \left[\log \frac{q_{\phi}(\mathbf{z}|\mathbf{x})}{p(\mathbf{z})}\right]}_{\text{KL Divergence}}
    %&= \mathbb{E}_{\mathbf{z}\sim q_{\theta}(\mathbf{z}|\mathbf{x})} \left[ \log \left[ \frac{p_{\theta}(\mathbf{x},\mathbf{z})}{q_{\phi}(\mathbf{z} | \mathbf{x})} \frac{q_{\phi}(\mathbf{z}|\mathbf{x})}{p_{\theta}(\mathbf{z}|\mathbf{x})} \right] \right] \\
    %&= \underbrace{\mathbb{E}_{\mathbf{z}\sim q_{\phi}(\mathbf{z}|\mathbf{x})} \left[ \log \frac{p_{\theta}(\mathbf{x}, \mathbf{z})}{q_{\phi}(\mathbf{z}|\mathbf{x})} \right]}_{\text{ELBO Term}} + \underbrace{\mathbb{E}_{\mathbf{z}\sim q_{\phi}(\mathbf{z}|\mathbf{x})} \left[ \log \frac{q_{\phi}(\mathbf{z}|\mathbf{x})}{p_{\theta}(\mathbf{z}|\mathbf{x})} \right]}_{\text{KL Divergence}}.
\end{align}
We can see that for any data point $\mathbf{x}$, the log-likelihood satisfies: 
\begin{align}
    \log p_{\theta}(\mathbf{x}) &\geq \mathcal{L}_{ELBO}(\theta, \phi; \mathbf{x}) \\
    & = \mathbb{E}_{\mathbf{z}\sim q_{\phi}(\mathbf{z}|\mathbf{x})} \left[\log p_{\theta}(\mathbf{x}|\mathbf{z})\right] - \mathcal{D}_{KL}(q_{\phi}(\mathbf{z}|\mathbf{x}) || p(\mathbf{z})).
\end{align}
%The first term is the variational lower bound (ELBO), which can be expressed as:
%\begin{equation}\label{eq:elbos}
%    \mathcal{L}_{ELBO}(\theta, \phi; \mathbf{x}) = \mathbb{E}_{\mathbf{z}\sim q_{\phi}(\mathbf{z}|\mathbf{x})} \left[ \log p_{\theta}(\mathbf{x},\mathbf{z}) - \log q_{\phi}(\mathbf{z}|\mathbf{x}) \right]
%\end{equation}
%The second term is the KL Divergence between the approximate posterior and the true posterior, $D_{KL}(q_{\phi}(\mathbf{z}|\mathbf{x}) || p_{\theta}(\mathbf{z}|\mathbf{x}))$, which is always non-negative. 
%Because of this, eq. (\ref{eq:elbos}) is a lower bound of the log-likelihood of the data point $\mathbf{x}$: 
%\begin{equation}
%    \mathcal{L}_{ELBO}(\theta, \phi; \mathbf{x}) = \log p_{\theta}(\mathbf{x}) - D_{KL}(q_{\phi}(\mathbf{z}|\mathbf{x}) || p_{\theta}(\mathbf{z}|\mathbf{x})) \leq \log p_{\theta}(\mathbf{x}).
%\end{equation}
If we maximize the previous loss with respect to the parameters $\theta$ and $\phi$, we are optimizing: 
\begin{itemize}
    \item \textbf{Reconstruction}: the first term encourage to maximize the likelihood $p_{\theta}(\mathbf{x}|\mathbf{z})$, which corresponds to the reconstruction of the data point $\mathbf{x}$ from the latent variables $\mathbf{z}$. Optimizing this term alone risks memorizing the training data, so we need extra regularization. 
    \item \textbf{Latent KL Regularization}: the second term encourages the encoder distribution $q_{\phi}(\mathbf{z}|\mathbf{x})$ to be close to the prior distribution $p(\mathbf{z})$, which is easy to sample from. 
\end{itemize}
Once we have the definition of the training objective, we can use stochastic gradient descent (SGD) to perform joint optimization w.r.t. all parameters $\theta$ and $\phi$. 
However, we have a problem with the first term of the ELBO because we can not backpropagate through the sampling process of $\mathbf{z}\sim q_{\phi}(\mathbf{z}|\mathbf{x})$, 
as we can see in the left part of Figure \ref{fig:reparameterization}. To solve this problem we can use the \textbf{reparameterization trick} which consists in expressing the sampling 
process as a deterministic function of the parameters and some noise. For example, if we assume that $q_{\phi}(\mathbf{z}|\mathbf{x})$ is a Gaussian distribution with mean $\mu$ and 
variance $\sigma^2$, we can sample from it as follows: 
\begin{equation}
    \mathbf{z} = \mu + \sigma \odot \epsilon, \quad \epsilon \sim \mathcal{N}(0, I).
\end{equation}
This way, we can backpropagate through the deterministic function and optimize the parameters $\phi$ of the encoder. This is shown in the right part of Figure \ref{fig:reparameterization}.

\begin{figure}   
\centering
\begin{tikzpicture}[
    >=Stealth,
    node distance=2.0cm, 
    % Scaled-down styles for nodes
    block/.style={draw=blue!70!black, fill=blue!5, rectangle, minimum height=0.9cm, minimum width=2.6cm, align=center, rounded corners=2pt, thick, font=\footnotesize},
    stoch_node/.style={draw=red!70!black, fill=red!5, circle, minimum size=1.5cm, align=center, thick, text width=1.3cm, font=\scriptsize},
    math_op/.style={draw=green!70!black, fill=green!5, diamond, minimum size=1.6cm, align=center, thick, inner sep=0pt, font=\footnotesize},
    io/.style={font=\footnotesize\bfseries},
    % Arrow styles
    forward_arrow/.style={->, thick, draw=black!80},
    grad_fail/.style={->, thick, dashed, draw=red!70!black},
    grad_success/.style={->, thick, dashed, draw=green!60!black}
]

    % ==================== Flow 1: Direct Sampling (LEFT COLUMN) ====================
    % Center of left column is at x=0
    \node (title1) at (0, 0) [font=\footnotesize\bfseries] {1. Direct Sampling};

    \node (encoder1) [block, below=0.3cm of title1] {Encoder\\$x \rightarrow \mu, \sigma^2$};
    \node (mu_sigma1) [below=0.3cm of encoder1, font=\footnotesize] {$\mu, \sigma^2$};
    \node (sampling1) [stoch_node, below=0.3cm of mu_sigma1] {Stochastic\\Sampling};
    \node (z1) [io, below=0.3cm of sampling1] {$z$};
    \node (decoder1) [block, below=0.3cm of z1] {Decoder\\$z \rightarrow \hat{x}$};

    % Forward Arrows (Left)
    \draw [forward_arrow] (encoder1) -- (mu_sigma1);
    \draw [forward_arrow] (mu_sigma1) -- (sampling1);
    \draw [forward_arrow] (sampling1) -- (z1);
    \draw [forward_arrow] (z1) -- (decoder1);

    % Compact Gradient Bus (Routed cleanly on the left)
    \coordinate (grad_bus1) at ($(decoder1.west) + (-1.0cm, 0)$);
    
    \draw [grad_fail] (decoder1.west) -- (grad_bus1) node[midway, above, font=\tiny\bfseries, text=red!70!black, inner sep=1pt] {Backprop} -- (grad_bus1 |- sampling1.west) -- (sampling1.west);
    \draw [grad_fail] (grad_bus1 |- z1.west) -- (z1.west);
    
    % Broken gradient indicator (Fixed alignment and overlap)
    \draw [grad_fail] (sampling1.north west) -- ++(-0.4cm,0.4cm) node[above left, font=\tiny\bfseries, text=red!70!black, inner sep=1pt, xshift=0.1cm, yshift=-0.1cm] {Broken};
    \coordinate (break_pt1) at ($(sampling1.north west) + (-0.2cm, 0.2cm)$);
    \draw [thick, red!70!black] ($(break_pt1)+(-0.1cm,-0.1cm)$) -- +(0.2cm,0.2cm);
    \draw [thick, red!70!black] ($(break_pt1)+(-0.1cm,0.1cm)$) -- +(0.2cm,-0.2cm);


    % ==================== Flow 2: Reparameterization (RIGHT COLUMN) ====================
    % Center of right column is slightly wider to give the green gradient bus room
    \node (title2) at (5.4cm, 0) [font=\footnotesize\bfseries] {2. Reparameterization};

    \node (encoder2) [block, below=0.3cm of title2] {Encoder\\$x \rightarrow \mu, \sigma^2$};
    \node (mu_sigma2) [below=0.3cm of encoder2, font=\footnotesize] {$\mu, \sigma^2$};
    
    \node (math_op2) [math_op, below=0.3cm of mu_sigma2] {$ \mu \!+\! \sigma \epsilon $};
    
    % Noise Node
    \node (epsilon2) [stoch_node, right=0.5cm of math_op2] {Noise $\epsilon$\\ \tiny $\mathcal{N}(0, I)$};

    \node (z2) [io, below=0.3cm of math_op2] {$z$};
    \node (decoder2) [block, below=0.3cm of z2] {Decoder\\$z \rightarrow \hat{x}$};

    % Forward Arrows (Right)
    \draw [forward_arrow] (encoder2) -- (mu_sigma2);
    \draw [forward_arrow] (mu_sigma2) -- (math_op2);
    
    % FIXED: Split the bidirectional arrows between Diamond and Noise so they don't overlap
    % 1. Forward stochastic path (Shifted slightly up)
    \draw [forward_arrow] ([yshift=3pt]epsilon2.west) -- ([yshift=3pt]math_op2.east) node[midway, above, font=\tiny, inner sep=1pt] {Stoch.};
    
    % FIXED: "Det." text moved slightly right so it doesn't clip the diamond
    \draw [forward_arrow] (math_op2.south) -- (z2.north) node[midway, right, font=\tiny, inner sep=2pt] {Det.};
    \draw [forward_arrow] (z2) -- (decoder2);

    % Compact Gradient Bus (Right) - Now sitting comfortably on the left of the blocks
    \coordinate (grad_bus2) at ($(decoder2.west) + (-1.0cm, 0)$);
    \coordinate (grad_bus2_top) at ($(encoder2.west) + (-1.0cm, 0)$);
    
    \draw [grad_success] (decoder2.west) -- (grad_bus2) node[midway, above, font=\tiny\bfseries, text=green!60!black, inner sep=1pt] {Backprop} -- (grad_bus2_top) -- (encoder2.west);
    
    % Connecting intermediate nodes to success bus
    \draw [grad_success] (grad_bus2 |- z2.west) -- (z2.west);
    \draw [grad_success] (grad_bus2 |- math_op2.west) -- (math_op2.west);
    \draw [grad_success] (grad_bus2 |- mu_sigma2.west) -- (mu_sigma2.west);
    
    % FIXED: Backward gradient fail path (Shifted slightly down to separate from forward path)
    \draw [grad_fail] ([yshift=-3pt]math_op2.east) -- ([yshift=-3pt]epsilon2.west) node[midway, below, font=\tiny\bfseries, text=red!70!black, inner sep=1pt] {No grad};
    
    % Cross mark correctly positioned directly on the shifted "No grad" line
    \coordinate (cross_pt) at ($([yshift=-3pt]math_op2.east)!0.5!([yshift=-3pt]epsilon2.west)$);
    \draw [thick, red!70!black] ($(cross_pt)+(-0.1cm,-0.1cm)$) -- +(0.2cm,0.2cm);
    \draw [thick, red!70!black] ($(cross_pt)+(-0.1cm,0.1cm)$) -- +(0.2cm,-0.2cm);

    % Divider line perfectly centered between the two column bounding boxes
    \draw [dashed, draw=gray!40, thick] (2.7cm, 0.2cm) -- (2.7cm, -7.2cm);

\end{tikzpicture}
\caption{Reparameterization Trick in VAEs}
\label{fig:reparameterization}
\end{figure}


% ******
% Gaussian VAEs 
% ******
\section{Gaussian VAEs}
The common choice for both the encoder and the decoder distributions in VAEs is the Gaussian distribution. This is because it allows us to use the reparameterization trick and 
also because it has a closed-form expression for the KL divergence. \\
The \textbf{encoder} $q_{\phi}(\mathbf{z}|\mathbf{x})$ is modelled as a Gaussian distribution defined as:
\begin{equation}
    q_{\phi}(\mathbf{z}|\mathbf{x}) = \mathcal{N} \left(\mathbf{z}; \mu_{\phi}(\mathbf{x}), \text{diag}(\sigma^2_{\phi}(\mathbf{x})) \right),
\end{equation}
where $\mu_{\phi}(\mathbf{x})$ and $\sigma^2_{\phi}(\mathbf{x})$ are the mean and variance of the Gaussian distribution, which are parameterized by a neural network with parameters $\phi$. \\
The \textbf{decoder} $p_{\theta}(\mathbf{x}|\mathbf{z})$ is also modelled as a Gaussian distribution defined as:
\begin{equation}
    p_{\theta}(\mathbf{x}|\mathbf{z}) = \mathcal{N} \left(\mathbf{x}; \mu_{\theta}(\mathbf{z}), \sigma^2\mathbf{I} \right),
\end{equation}
where $\mu_{\theta}(\mathbf{z})$ is the mean of the Gaussian distribution, which is parameterized by a neural network with parameters $\theta$, and $\sigma^2\mathbf{I}$ is the covariance matrix, 
which is usually fixed to be a diagonal matrix with constant variance. \\ 
\\
Under this Gaussian assumption, we can obtain a closed-form expression for the ELBO loss function. The reconstruction term can be expressed as:
\begin{equation}
    \mathbb{E}_{\mathbf{z}\sim q_{\phi}(\mathbf{z}|\mathbf{x})} \left[\log p_{\theta}(\mathbf{x}|\mathbf{z})\right] = -\frac{1}{2\sigma^2} \mathbb{E}_{\mathbf{z}\sim q_{\phi}(\mathbf{z}|\mathbf{x})} \left[||\mathbf{x} - \mu_{\theta}(\mathbf{z})||^2\right] + C,
\end{equation}
where $C$ is a constant that does not depend on the parameters. The KL divergence term can be expressed as:
\begin{equation}
    \mathcal{D}_{KL}(q_{\phi}(\mathbf{z}|\mathbf{x}) || p(\mathbf{z})) = \frac{1}{2} \sum_{i=1}^k \left( \sigma^2_{\phi,i}(\mathbf{x}) + \mu^2_{\phi,i}(\mathbf{x}) - 1 - \log \sigma^2_{\phi,i}(\mathbf{x}) \right),
\end{equation}
where $k$ is the dimension of the latent space. Therefore, the ELBO loss function can be expressed as:
\begin{align}
    \mathcal{L}_{ELBO}(\theta, \phi; \mathbf{x}) &= -\frac{1}{2\sigma^2} \mathbb{E}_{\mathbf{z}\sim q_{\phi}(\mathbf{z}|\mathbf{x})} \left[||\mathbf{x} - \mu_{\theta}(\mathbf{z})||^2\right] \\ 
    &- \frac{1}{2} \sum_{i=1}^k \left( \sigma^2_{\phi,i}(\mathbf{x}) + \mu^2_{\phi,i}(\mathbf{x}) - 1 - \log \sigma^2_{\phi,i}(\mathbf{x}) \right) + C.
\end{align} 
Despite their mathematical elegance and continuous latent spaces, standard Gaussian Variational Autoencoders notoriously suffer from generating blurry images. 
This limitation fundamentally stems from the VAE objective function and its underlying statistical assumptions. The problem can be attributed to four primary factors:
\begin{itemize}
    \item \textbf{MSE Loss}: In a standard Gaussian VAE with a fixed variance $\sigma^2 \mathbf{I}$, maximizing the likelihood $p_\theta(\mathbf{x}|\mathbf{z})$ is mathematically equivalent to 
    minimizing the $\mathcal{L}_2$ distance between the input $\mathbf{x}$ and the reconstruction $\hat{\mathbf{x}}$:
    \begin{equation}
        \log p_\theta(\mathbf{x}|\mathbf{z}) \propto -||\mathbf{x} - \mu_\theta(\mathbf{z})||^2_2
        \label{eq:vae_mse}
    \end{equation}
    The $\mathcal{L}_2$ loss heavily penalizes large pixel-wise deviations. When the model is uncertain about the exact location of high-frequency details (such as sharp edges), it outputs the statistical 
    mean of all plausible variations to minimize the overall penalty. This phenomenon, known as \textit{mode averaging}, results in safe, blurry predictions rather than rendering sharp, distinct features.
    \item \textbf{Information Bottleneck}: The ELBO includes a regularization term that minimizes the KL divergence between the approximate posterior and the prior:
    \begin{equation}
        \mathcal{D}_{\text{KL}}(q_\phi(\mathbf{z}|\mathbf{x}) \parallel p(\mathbf{z}))
        \label{eq:vae_kl}
    \end{equation}
    This term acts as an information bottleneck, forcing the encoded latent representations to closely match the simple prior, typically $\mathcal{N}(\mathbf{0}, \mathbf{I})$. Consequently, the encoder 
    discards complex, high-frequency spatial features that are difficult to compress, prioritizing global structures and low-frequency components instead.
    \item \textbf{Posterior Collapse}: Furthermore, consistent with findings in \cite{bowman2016generatingsentencescontinuousspace} and \cite{sonderby2016howtrain}, stochastic optimization using the unmodified lower bound objective frequently 
    gets stuck in an undesirable stable equilibrium. At the onset of training, the reconstruction likelihood term $\log p_\theta(\mathbf{x}|\mathbf{z})$ is relatively weak. Consequently, an initially 
    attractive state for the network is to aggressively minimize the latent cost, driving the approximate posterior to match the prior, $q_\phi(\mathbf{z}|\mathbf{x}) \approx p(\mathbf{z})$. This results 
    in a stable equilibrium---often termed \textit{posterior collapse}---from which the model struggles to escape, yielding uninformative latent representations and contributing to poor generation quality. 
    To mitigate this, a common solution proposed in \cite{bowman2016generatingsentencescontinuousspace, sonderby2016howtrain} is to implement an optimization schedule (KL annealing) where the weight of the latent cost 
    $\mathcal{D}_{\text{KL}}(q_\phi(\mathbf{z}|\mathbf{x}) \parallel p(\mathbf{z}))$ is slowly annealed from 0 to 1 over multiple epochs.
\end{itemize}


% ******
% DDPMs 
% ******
\section{Denoising Diffusion Probabilistic Models}
One important technique developed to enhance the performance of VAEs and overcome the main limitations was \textbf{Hierarchical Variational Autoencoders} (HVAEs) \cite{vahdat2021nvaedeephierarchicalvariational} which are a type of VAE that uses a hierarchical 
structure in the latent space. This hierarchical structure allows the model to capture more complex and high-frequency details in the data, which can help to reduce the blurriness of the generated images. A visualization of the 
architecture of HVAEs is shown in Figure \ref{fig:hvvae}. In this architecture, we have $L$ layers of latent variables $\mathbf{z}_1, \mathbf{z}_2, \dots, \mathbf{z}_L$, where each layer captures different levels of abstraction. \\
In this setup, we can define the following factorization of the joint distribution: 
\begin{figure}   
\centering
\begin{tikzpicture}[
    % Tighter styles for narrower fit
    block/.style={
        draw,
        thick,
        rectangle,
        rounded corners=1mm,
        minimum width=1.1cm,
        minimum height=1.8cm,
        fill=black!5,
        font=\large,
        align=center
    },
    arr/.style={->, thick, >=Stealth},
    revarr/.style={<-, thick, >=Stealth},
    % Added a label style to make the math font slightly smaller and reduce padding
    lbl/.style={midway, font=\footnotesize, inner sep=5.5pt} 
]

    % --- Nodes ---
    \node[block] (x) {$\mathbf{x}$};
    \node[block, right=1.5cm of x] (z1) {$\mathbf{z}_1$};
    \node[block, right=1.5cm of z1] (z2) {$\mathbf{z}_2$};
    
    % Reduced the space around the ellipsis
    \node[right=0.7cm of z2, font=\Large, inner xsep=0mm] (dots) {$\dots$};
    
    \node[block, right=1.8cm of dots] (zl) {$\mathbf{z}_L$};

    % --- Edges ---
    % Brought the arrows slightly closer together (yshift=3.5mm instead of 5mm)
    
    % Connections between x and z1
    \draw[arr]    ([yshift=3.5mm]x.east) -- ([yshift=3.5mm]z1.west)
        node[lbl, above] {$q_\theta(\mathbf{z}_1|\mathbf{x})$};
    \draw[revarr] ([yshift=-3.5mm]x.east) -- ([yshift=-3.5mm]z1.west)
        node[lbl, below] {$p_\phi(\mathbf{x}|\mathbf{z}_1)$};

    % Connections between z1 and z2
    \draw[arr]    ([yshift=3.5mm]z1.east) -- ([yshift=3.5mm]z2.west)
        node[lbl, above] {$q_\theta(\mathbf{z}_2|\mathbf{z}_1)$};
    \draw[revarr] ([yshift=-3.5mm]z1.east) -- ([yshift=-3.5mm]z2.west)
        node[lbl, below] {$p_\phi(\mathbf{z}_1|\mathbf{z}_2)$};

    % Connections entering the ellipsis from z2
    \draw[arr]    ([yshift=3.5mm]z2.east) -- ([yshift=3.5mm]dots.west);
    \draw[revarr] ([yshift=-3.5mm]z2.east) -- ([yshift=-3.5mm]dots.west);

    % Connections leaving the ellipsis to zL
    \draw[arr]    ([yshift=3.5mm]dots.east) -- ([yshift=3.5mm]zl.west)
        node[lbl, above] {$q_\theta(\mathbf{z}_L|\mathbf{z}_{L-1})$};
    \draw[revarr] ([yshift=-3.5mm]dots.east) -- ([yshift=-3.5mm]zl.west)
        node[lbl, below] {$p_\phi(\mathbf{z}_{L-1}|\mathbf{z}_L)$};

\end{tikzpicture}
\caption{Hierarchical VAE architecture with $L$ layers of latent variables. Each layer captures different levels of abstraction, allowing for richer representations and improved generation quality.}
\label{fig:hvvae}
\end{figure}
\begin{equation}
    p_{\theta}(\mathbf{x},\mathbf{z}_{1:L}) = p_{\theta}(\mathbf{x}|\mathbf{z}_1)\prod_{i=2}^{L} p_{\theta}(\mathbf{z}_{i-1}|\mathbf{z}_i)p(\mathbf{z}_L),
\end{equation}
and the marginal distribution of the data point $\mathbf{x}$ can be expressed as:
\begin{equation}
    p_{\theta}(\mathbf{x}) = \int p_{\theta}(\mathbf{x},\mathbf{z}_{1:L})d\mathbf{z}_{1:L},
\end{equation}
where the generation is performed progressively from the latent variable $\mathbf{z}_L$, which follows the prio distribution $p(\mathbf{z}_L)$, to the data point $\mathbf{x}$, which is generated from the latent variable $\mathbf{z}_1$.
The inference model can be defined as: 
\begin{equation}
    q_{\phi}(\mathbf{z}_{1:L}|\mathbf{x}) = q_{\phi}(\mathbf{z}_1|\mathbf{x})\prod_{i=2}^{L} q_{\phi}(\mathbf{z}_i|\mathbf{z}_{i-1}),
\end{equation}
where the inference is performed progressively from the data point $\mathbf{x}$ to the latent variable $\mathbf{z}_L$. \\
The ELBO loss is very similar to the one of the standard VAE, but we have to consider all the layers of latent variables:
\begin{align}
    \mathcal{L}_{ELBO}(\theta, \phi; \mathbf{x}) &= \mathbb{E}_{\mathbf{z}_{1:L}\sim q_{\phi}(\mathbf{z}_{1:L}|\mathbf{x})} \left[\log \frac{p_{\theta}(\mathbf{x},\mathbf{z}_{1:L})}{q_{\phi}(\mathbf{z}_{1:L}|\mathbf{x})} \right],
\end{align}
which can also be decomposed into the reconstruction term and the KL divergence term. \\
\\
This method of deep, stacked layers revolutionized the field of generative modeling and is the basis for many state-of-the-art models. The training of HVAEs present several challenges because the encoder 
and decoder must be trained jointly and learning becomes unstable. A clever twist to solve the main problems and maintain the hierarchical structure is the \textbf{Denoising Diffusion Probabilistic Models} (DDPMs)
\cite{sohldickstein2015deepunsupervisedlearningusing, ho2020denoisingdiffusionprobabilisticmodels}, the next step in generative modeling, that can be defined with the following stochastic processes. \\
\\
\textbf{Forward Pass}. This is a fixed encoder (not learnable) that progressively adds noise to the data point $\mathbf{x}$ via a transition 
kernel $p(\mathbf{x}_t|\mathbf{x}_{t-1})$ until we get pure noise $\mathbf{x}_T \sim \mathcal{N}(\mathbf{0}, \mathbf{I})$. The fixed Gaussian transition kernel is defined as:
\begin{equation}\label{eq:forward_kernel}
    p(\mathbf{x}_t|\mathbf{x}_{t-1}) = \mathcal{N}(\mathbf{x}_t; \sqrt{1-\beta_t}\mathbf{x}_{t-1}, \beta_t\mathbf{I}),
\end{equation} 
where $\mathbf{x}_0\sim p_{data}$, $\mathbf{x}_T \sim \mathcal{N}(\mathbf{0}, \mathbf{I})$ and the sequence $\{\beta_t\}_{t=1}^T\in(0,1)$ is the 
noise schedule, which controls the amount of noise added at each step, as we can clearly see in the following formulation: 
\begin{equation}\label{eq:forward_sample}
    \mathbf{x}_t = \sqrt{\alpha_t} \mathbf{x}_{t-1} + \sqrt{1-\alpha_t} \epsilon, \quad \epsilon \sim \mathcal{N}(\mathbf{0}, \mathbf{I}),
\end{equation}
where $\alpha_t = 1-\beta_t$. Another important formulation of the forward pass can be achieve by recursively applying the transition kernel (Equation \ref{eq:forward_kernel}), 
where we achieve a closed-form expression for the distribution of noisy samples at step $t$ given the original data:
\begin{equation}
    p(\mathbf{x}_t|\mathbf{x}_0) = \mathcal{N}(\mathbf{x}_t; \sqrt{\bar{\alpha}_t}\mathbf{x}_0, (1-\bar{\alpha}_t)\mathbf{I}),
\end{equation}
where $\bar{\alpha}_t = \prod_{s=1}^T \alpha_s$. This formulation is very useful because it allows us to sample from the forward process at any 
time step $t$ without having to recursively apply the transition kernel $t$ times. We can also define the sampling of $\mathbf{x}_t$ as in Equation \ref{eq:forward_sample} as follows:
\begin{equation}\label{eq:forward_sample2}
    \mathbf{x}_t = \sqrt{\bar\alpha_t} \mathbf{x}_0 + \sqrt{1-\bar{\alpha}_t} \epsilon, \quad \epsilon \sim \mathcal{N}(\mathbf{0}, \mathbf{I}).
\end{equation}
\\
\textbf{Reverse Pass}: This is the decoder, where a neural network is trained to reverse the noising process by learning the parameterized 
distribution $p_{\theta}(\mathbf{x}_{t-1}|\mathbf{x}_t)$. If we start from a pure noise sample, we can iteratively apply the reverse process 
to generate a data point $\mathbf{x}_0$ that resembles the training data, following a Markov chain. The training goal is to approximate the unknown reverse transition 
kernel $p(\mathbf{x}_{t-1}|\mathbf{x}_t)$ with a learnable parametric model $p_{\theta}(\mathbf{x}_{t-1}|\mathbf{x}_t)$, using the KL divergence 
as the training objective: 
\begin{equation}\label{eq:kl_reverse}
    \mathbb{E}_{\mathbf{x}_t \sim p_t} \left[ \mathcal{D}_{KL}(p(\mathbf{x}_{t-1}|\mathbf{x}_t) || p_{\theta}(\mathbf{x}_{t-1}|\mathbf{x}_t)) \right].
\end{equation}
We can express the unknown reverse transition kernel $p(\mathbf{x}_{t-1}|\mathbf{x}_t)$ in terms of the forward process using Bayes' rule:
\begin{equation}
    p(\mathbf{x}_{t-1}|\mathbf{x}_t) = \frac{p(\mathbf{x}_t|\mathbf{x}_{t-1}) p_{t-1}(\mathbf{x}_{t-1})}{p_t(\mathbf{x}_t)},
\end{equation}
where the marginal distribution $p_t(\mathbf{x}_t)$ can be expressed as:
\begin{equation}
    p_t(\mathbf{x}_t) = \int p(\mathbf{x}_t|\mathbf{x}_0) p_{data}(\mathbf{x}_0) d\mathbf{x}_0.
\end{equation}
Since $p_{data}$ is unknown, we can not compute the marginal distribution $p_t(\mathbf{x}_t)$ thus the computation of the target distribution 
$p(\mathbf{x}_{t-1}|\mathbf{x}_t)$ is intractable. A trick to solve this problem is to condition the reverse transition kernel on the original data 
point $\mathbf{x}_0$. Now we can define the target distribution as follows:
\begin{equation}
    p(\mathbf{x}_{t-1}|\mathbf{x}_t, \mathbf{x}_0) = \frac{p(\mathbf{x}_t|\mathbf{x}_{t-1}) p_{t-1}(\mathbf{x}_{t-1}|\mathbf{x}_0)}{p_t(\mathbf{x}_t|\mathbf{x}_0)},
\end{equation}
which allows us to have a tractable expression and to derive a tractable objective equivalent to the previous intractable KL divergence in Equation \ref{eq:kl_reverse}. 
As a result, $p(\mathbf{x}_{t-1}|\mathbf{x}_t, \mathbf{x}_0)$ is a Gaussian distribution with mean and variance that can be computed in closed-form: 
\begin{equation}
    p(\mathbf{x}_{t-1}|\mathbf{x}_t, \mathbf{x}_0) = \mathcal{N}(\mathbf{x}_{t-1}; \tilde{\mu}_t(\mathbf{x}_t, \mathbf{x}_0, t), \sigma^2(t)\mathbf{I}),
\end{equation}
where 
\begin{equation}\label{eq:tilde_mean}
    \tilde{\mu}_t(\mathbf{x}_t, \mathbf{x}_0, t) = \frac{\sqrt{\bar{\alpha}_{t-1}}\beta_t}{1-\bar{\alpha}_t}\mathbf{x}_0 + \frac{(1-\bar{\alpha}_{t-1})\sqrt{\alpha}_t}{1-\bar{\alpha}_t}\mathbf{x}_t, \quad \sigma^2(t) = \frac{1-\bar{\alpha}_{t-1}}{1-\bar{\alpha}_t}\beta_t.
\end{equation}
Finally, DDPM assumes that each reverse transition kernel $p_{\theta}(\mathbf{x}_{t-1}|\mathbf{x}_t)$ is a Gaussian distribution with mean $\mu_{\theta}(\mathbf{x}_t, t)$ 
and variance $\sigma^2(t)\mathbf{I}$, where the variance is fixed and the mean is parameterized by a neural network with parameters $\theta$. It can be defined as 
\begin{equation}
    p_{\theta}(\mathbf{x}_{t-1}|\mathbf{x}_t) = \mathcal{N}(\mathbf{x}_{t-1}; \mu_{\theta}(\mathbf{x}_t, t), \sigma^2(t)\mathbf{I}).
\end{equation}
Using the previous conditional derivation of the unknown reverse transition kernel, we can define a closed-form expression for the KL divergence training loss:
\begin{equation}\label{eq:loss_ddpm}
    \mathcal{L}(\mathbf{x}_0;\theta) = \mathbb{E}_{t, \mathbf{x}_0} \left[\frac{1}{2\sigma^2(t)} ||\mu_{\theta}(\mathbf{x}_t, t) - \tilde{\mu}_t(\mathbf{x}_t, \mathbf{x}_0, t)||^2 + C \right].
\end{equation}
where $C$ is a constant that does not depend on the parameters. This loss can be interpreted as a weighted $\mathcal{L}_2$ loss between the predicted mean 
$\mu_{\theta}(\mathbf{x}_t, t)$ and the true mean $\tilde{\mu}_t(\mathbf{x}_t, \mathbf{x}_0, t)$ of the reverse transition kernel. \\
In practice, we do not use the loss of Equation \ref{eq:loss_ddpm} directly. If we combine the definition of the true mean $\tilde{\mu}_t(\mathbf{x}_t, \mathbf{x}_0, t)$ in Equation 
\ref{eq:tilde_mean} with Equation \ref{eq:forward_sample2}, we can obtain two different parameterizations of the mean that can be used to derive two different but equivalent training 
losses that are easier to optimize. 
\begin{itemize}
    \item $\epsilon$-Prediction. From Equation \ref{eq:forward_sample2} we obtain $\mathbf{x}_0 = \frac{1}{\sqrt{\bar{\alpha}_t}}(\mathbf{x}_t - \sqrt{1-\bar{\alpha}_t}\epsilon)$ and we can 
    substitute this expression in the definition of the true mean $\tilde{\mu}_t(\mathbf{x}_t, \mathbf{x}_0, t)$ to obtain the following expression:
    \begin{equation}
        \tilde{\mu}_t(\mathbf{x}_t, \mathbf{x}_0, t) = \frac{1}{\sqrt{\alpha}_t} \left( \mathbf{x}_t - \frac{\beta_t}{\sqrt{1-\bar{\alpha}_t}} \epsilon \right),
    \end{equation}
    which can also be applied to parameterized mean using now a neural network that predicts the noise $\epsilon_{\theta}(\mathbf{x}_t, t)$ instead of the mean $\mu_{\theta}(\mathbf{x}_t, t)$:
    \begin{equation}
    \mu_{\theta}(\mathbf{x}_t, t) = \frac{1}{\sqrt{\alpha}_t} \left( \mathbf{x}_t - \frac{\beta_t}{\sqrt{1-\bar{\alpha}_t}} \epsilon_{\theta}(\mathbf{x}_t, t) \right).
    \end{equation}
    We can define a new loss function based on the noise prediction as follows:
    \begin{equation}
        \mathcal{L}(\mathbf{x}_0;\theta) = \mathbb{E}_{t, \mathbf{x}_0, \epsilon} \left[\frac{\beta_t^2}{2\sigma^2(t)\alpha_t(1-\bar{\alpha}_t)} ||\epsilon_{\theta}(\mathbf{x}_t, t) - \epsilon||^2 \right],
    \end{equation}
    where a simple version of this loss, where the weighting factor is removed, is commonly used in practice for the good performance and stability of the training process found in \cite{ho2020denoisingdiffusionprobabilisticmodels}.
    The final loss can be expressed as: 
    \begin{equation}
        \mathcal{L}(\mathbf{x}_0;\theta) = \mathbb{E}_{t, \mathbf{x}_0, \epsilon} \left[||\epsilon_{\theta}(\mathbf{x}_t, t) - \epsilon||^2 \right].
    \end{equation}
    \item $\mathbf{x}_0$-Prediction: With the definition of the true mean $\tilde{\mu}_t(\mathbf{x}_t, \mathbf{x}_0, t)$ (Equation \ref{eq:tilde_mean}) we can obtain 
    a new parameterization of the mean by using a neural network that predicts the original data point $\mathbf{x}_0$ instead of the mean $\mu_{\theta}(\mathbf{x}_t, t)$:
    \begin{equation}
        \mu_{\theta}(\mathbf{x}_t, t) = \frac{\sqrt{\bar{\alpha}_{t-1}}\beta_t}{1-\bar{\alpha}_t}\mathbf{x}_{\theta}(\mathbf{x}_t, t) + \frac{(1-\bar{\alpha}_{t-1})\sqrt{\alpha}_t}{1-\bar{\alpha}_t}\mathbf{x}_t,
    \end{equation}
    which, as in the previous case, can be used to derive a new loss function based on the original data point prediction as follows:
    \begin{equation}
        \mathcal{L}(\mathbf{x}_0;\theta) = \mathbb{E}_{t, \mathbf{x}_0, \epsilon} \left[\omega_t||\mathbf{x}_{\theta}(\mathbf{x}_t, t) - \mathbf{x}_0||^2 \right].
    \end{equation}
    where $\omega_t$ is a weighting factor that depends on the noise schedule and the variance of the reverse transition kernel, which can be ignored in practice. 
\end{itemize}








\newpage
\begin{summarybox}[label={box:latent_math}]{VAEs and DDPMs}
    This box summarizes the main mathematical equations of VAEs and DDPMs:

    \medskip
    \textbf{1. Training Objective of VAEs:}
    \begin{equation}
        \mathcal{L}_{ELBO}(\theta, \phi; \mathbf{x}) = \mathbb{E}_{\mathbf{z}\sim q_{\phi}(\mathbf{z}|\mathbf{x})} \left[\log p_{\theta}(\mathbf{x}|\mathbf{z})\right] - \mathcal{D}_{KL}(q_{\phi}(\mathbf{z}|\mathbf{x}) || p(\mathbf{z})).
    \end{equation}

    \medskip
    \textbf{2. Training Objective of Gaussian VAEs:}
    \begin{align}
        &\mathcal{L}_{ELBO}(\theta, \phi; \mathbf{x}) = -\frac{1}{2\sigma^2} \mathbb{E}_{\mathbf{z}\sim q_{\phi}(\mathbf{z}|\mathbf{x})} \left[||\mathbf{x} - \mu_{\theta}(\mathbf{z})||^2\right] \\ 
        &- \frac{1}{2} \sum_{i=1}^k \left( \sigma^2_{\phi,i}(\mathbf{x}) + \mu^2_{\phi,i}(\mathbf{x}) - 1 - \log \sigma^2_{\phi,i}(\mathbf{x}) \right) + C.
    \end{align} 

    \medskip
    \textbf{3. Forward Pass in DDPMs:}
    \begin{align}
        & p(\mathbf{x}_t|\mathbf{x}_{t-1}) = \mathcal{N}(\mathbf{x}_t; \sqrt{1-\beta_t}\mathbf{x}_{t-1}, \beta_t\mathbf{I}) \\
        & p(\mathbf{x}_t|\mathbf{x}_0) = \mathcal{N}(\mathbf{x}_t; \sqrt{\bar{\alpha}_t}\mathbf{x}_0, (1-\bar{\alpha}_t)\mathbf{I})
    \end{align} 

    \medskip
    \textbf{4. $\epsilon$-Prediction loss in DDPMs:}
    \begin{equation}
        \mathcal{L}(\theta) = \mathbb{E}_{t, \mathbf{x}_0, \epsilon} \left[||\epsilon_{\theta}(\mathbf{x}_t, t) - \epsilon||^2 \right].
    \end{equation} 

    \medskip
    \textbf{5. $\mathbf{x}_0$-Prediction loss in DDPMs:}
    \begin{equation}
        \mathcal{L}(\theta) = \mathbb{E}_{t, \mathbf{x}_0, \epsilon} \left[\omega_t||\mathbf{x}_{\theta}(\mathbf{x}_t, t) - \mathbf{x}_0||^2 \right].
    \end{equation} 
    
    %\vsub
    %{\small \color{boxaccent}\textbf{Note:} Ensure latent dimensions $z_i$ and $z_j$ maintain near-zero mutual information to preserve explainability.}
\end{summarybox} 


%*****************************************
%*****************************************
%*****************************************
%*****************************************
%*****************************************